\paragraph{从拟合参数到未来位置预测。}
每粒子从$x_0$出发，按$\pi$抽取模式；Full在整条路径上保持该模式。
对每个不超过5个保存秒的事件区间$\Delta$，将局部仿射漂移写为
$\mathsf A_mx+a_n$，在条件均值处求非位置特征并冻结。
保留仿射核及条件矩更新为
\begin{align}
 F&=e^{\Delta\mathsf A_m},\quad
 J=\int_0^\Delta e^{v\mathsf A_m}\,dv,\quad
 \Omega=\int_0^\Delta e^{v\mathsf A_m}Q_m e^{v\mathsf A_m^\top}\,dv,
 \label{eq:main-affine-kernel}\\
 \mu_{n+1}&=F\mu_n+Ja_n,\qquad
 C_{n+1}=FC_nF^\top+\Omega,\label{eq:main-moment-update}\\
 X_{n+1}^{(p)}&=FX_n^{(p)}+Ja_n+Gz_{n,p},\quad
 GG^\top=\Omega,\quad z_{n,p}\sim\mathcal N(0,I_2).
 \label{eq:main-particle-update}
\end{align}
初始$\mu_0=x_0$、$C_0=0$，同一带噪路径提供四个实际目标时刻的位置。
该核对冻结仿射子问题精确；在条件均值处求太阳特征是对完整非线性模型的近似。
保存程序中的“FP”因此指条件高斯矩传播，不是空间Fokker--Planck偏微分方程求解。
数值变体采用Euler--Maruyama或split核，MC／CRN变体沿个体路径求非线性条件。
完整定义见\supp{sec:method-model}。
该过程生成式\eqref{eq:forecast-object}的经验分布，评估预测没有另加半径校准步骤。
